\section{Adaptive Streaming}
While streaming a content we need to take into account rate-quality (tradeoff). Video is the most difficult problem because we have to reduce uncompressed video (from the order of Gbps of rate, down to some Mbps).\\
Possible applications of multimedia streaming:\begin{itemize}
    \item Bulk Transfer (need data integrity)
    \item VOD (delay tolerant)
    \item Live streaming (scalability and low delay)
\end{itemize}
It is possible to control the encoders such that we are able to compress video with the same rate.
\subsection{Video Content Distribution}
\subsubsection{Challenges}
Each user has different links to the central server, and each user has different constraints (display, battery,...)\\
A possible solution is the Conversational Paradigm (1 source to N destinations):
$$\boxed{\text{Speaker A}} \longrightarrow \boxed{\text{Conference SRV}} \longrightarrow \begin{cases}
    \boxed{\text{Listener B (TV)}}\\
    \boxed{\text{Listener C (Phone)}}\\
    \boxed{\text{Listener }n}
\end{cases}$$
How to choose the coding rate of my video?
\subsubsection{Coding Rate Bottleneck}
We need:
$$R_c \leq C_k \qquad (\text{capacity limit of user }k)$$
We need to keep the rate lower than the max capacity, to avoid buffering.\\
if $R_c = \min\limits_k c_k$ we fix a low resolution also for high-capacity linked users.\\
Solution: use MANE, also deployed in CDNs and 5G Networks.
\subsubsection{Paradigms}
\begin{minipage}{0.475\textwidth}
    \centering\textbf{Push}
    \begin{itemize}
        \item Goal: Minimize latency
        \item Stack: WebRTC/RTCP/RTP/UDP
        \item Server pushes, client listens\\
        (Real Time approach)
    \end{itemize}
\end{minipage}\hfill
\begin{minipage}{0.475\textwidth}
    \centering
    \textbf{Pull}
    \begin{itemize}
        \item Goal: Maximize QoE
        \item Stack:\\
        $[$DASH,HLS$]$/HTTP3/QUIC/UDP/(CDNs)
        \item Client asks for the content, server serves
    \end{itemize}
\end{minipage}
UDP is quick and fast, RTP just adds additional infos like:
\begin{itemize}
    \item Packet Numbering
    \item if Out-of-order packets $\to$ sort out
\end{itemize} 
RTCP also adds also a metrics exchange as control messages to let the server adapt the $R_c$.\\
There are also timestamps added to packets to handle sync and control jitter.
\subsection{Scalable Video Coding (SVC)}
Useful for a $1\to N$ realtime video streaming. IDEA: just perform 1 encoding and then be able to extract several versions of the content. (JP2K like capabilities $\to$ progressive video coding)
\subsubsection{Scalability Types}
\begin{itemize}
    \item \textbf{Time} Scalability: exploit temporal prediction
    \item \textbf{Quality} Scalability: use progressive coding, Base-Layer (BL) as prediction of ELs (Enhancement-Layers)
    \item \textbf{Space} Scalability: exploit Multiresolution Transforms, low-res base-layer + enhancement layers
\end{itemize}
Enhanced Layer is computed as $= I_{\text{nput}} - \text{DEC}(\text{ENC}( I_{\text{nput}}))$
\begin{center}
\includegraphics[width=0.5\linewidth]{svc_encoder.png}
\end{center}
\noindent SVC layered encoder: a coarse encoder $Q_B$ produces the \textbf{Base Layer}; its local decoder reconstruction is subtracted from the input and the residual is coded by $Q_E$ into the \textbf{Enhancement Layer} (quality scalability).
\subsubsection{Smart Layer Pruning}
Assume we have the following architecture:
$$\boxed{\text{Speaker}} \qquad \longrightarrow \qquad \boxed{\text{Confluence Server}} \qquad \longrightarrow \qquad \boxed{n\ U_i}$$
The speaker produces the Base-layer + 2 more enhancement layers, total of 3 Layers.\\
Confluence server decides how to send the layers to users, it's the \textbf{Selective Forward Unit}, usually a MANE.
\begin{center}
\includegraphics[width=\linewidth]{svc_conference.png}
\end{center}
\noindent SFU layer pruning: the speaker sends one unified bitstream (BL + EL1 + EL2). The conference server forwards \emph{all} layers to the high-capacity client (full quality) and only the Base Layer to the constrained mobile client --- zero re-encoding, just NAL-header inspection.
\subsubsection{Streaming via Legacy HTTP}
Why don't we use the same solution for streaming? SFU needs memory proportional to the connected users, to keep the state of each. And the BL + EL solution has some overhead w.r.t. a single high res block of data.
Solution: use HTTP because it is stateless and also largely implemented.\\
We can slice the video and have caching using the \textbf{CDNs}\footnote{Content Delivery Networks} at the network edges.\\
The only drawback of Legacy HTTP is TCP, we increase latency and stop playout if a segment is lost (Head-of-line blocking). This implies ReTX and large delays. Solution: QUIC.
\subsubsection{QUIC and HTTP/3}
HTTP/3 is based on UDP:
\begin{itemize}
    \item removes Head-of-line blocking
    \item it has short setup time (no handshakes)
    \item has connection migration, not bound to (ip,port) couple, just a connection ID
\end{itemize}
\subsection{Adaptive Bitrate Streaming}
\subsubsection{Media Presentation Description (MPD) file}
\begin{itemize}
    \item \textbf{Period}: It is representing a certain content
    \item \textbf{AdaptationSet}: It represents a certain track (video, audio,...)
    \item \textbf{Representation}: It represents the different qualities (different codecs, resolutions, bitrates,...)
    \item \textbf{Segments}: They are the actual sliced media parts, by joining them together we get the actual content.
\end{itemize}
Further notes:
\begin{itemize}
    \item Segments (.m4s) are Indepently decodable, they only need an init signature part (.m4i)
    \item To encode a segment you have to wait the IDR block in a certain GOP.\\
    Also in the decoding we'll start from a fresh new IDR block.
    \item If we know networks changes quickly: better short segments, otherwise viceversa.
    \item For Live streaming we have to use small chunks of about $200 - 500$ ms, this is achieved with CMAF\footnote{Common Media Application Format}\\
    this leads to a glass-to-glass latency of about $1 \to 3$ seconds.
\end{itemize}
\subsubsection{Network Metrics: Throughput}
Here we use continuous models to describe discrete things.
$$S(t) = \frac{d}{dt}D(t) \leq R_{\text{PHY}}$$
This inequality takes into account also for Overhead + Congestions + Other users on the link.
Since $S$ can be very sparky we want to use the average:
$$\bar{S} = \frac{1}{T}\int_0^TS(t)dt = \frac{D(T)}{T}$$
\subsubsection{Network Metrics: E-2-E Nodal Delay}
$$d_{e2e} = \sr{}{k \in \text{ links}}d_{\text{processing}}(k) + d_{\text{queue}}(k) + d_{\text{TX}}(k) + d_{\text{propagation}}(k)$$
\begin{itemize}
    \item Processing: Error control, header reading
    \item Queuing: buffering
    \item Transmission: from the 1st bit to the last one, $\frac{\text{pkt size}}{\text{link rate}}$
    \item Propagation: usually distance over speed of light.
\end{itemize}
In modern networks an increase of $R$ does affect only $d_{TX}$, the other delays keep the same size.\\
The only thing we can change is the distance, we use CDNs to bring the content nearer to the users.\\
Jitter is the most annoying phenomenon. Solution: use buffers to absorbe the variability of the delay (Jitter).
\subsection{Buffer Dynamics}
\subsubsection{Rebuffering event}
Rebuffering: process of filling our buffers, it involves the download of some segments to play later 
$$B(t) = L \cdot T_s \qquad \text{where }\begin{cases}
    L = \# \text{segments}\\
    T_s = \text{duration of a segment}
\end{cases}$$
If the network goes down, slowly buffer occupation decreases.\\
In a QoE maximization approach, for a user is better to have low quality video than an interruption.
\subsubsection{Playout management}
When I empty the buffer I wait for the playout to restart until I received $M$ segments (Hysteresis):
$$B(t) = M \cdot T_s \qquad \text{typically better for users if }L \geq M$$
%For users typically is better to have a $L \geq M$ contraint
We also set:
$$\frac{dB}{dt} = f_{IN} - f_{OUT}$$
Flows:
\begin{itemize}
    \item $f_{IN}$ Input Flow Rate, derived from network condition
        \begin{itemize}
            \item Average $f_{IN} = \frac{D(T)}{T} = \frac{\bar{S}}{R_c}$
            \item Instantaneous flow rate: $\frac{S(t)}{R_c} = \lim\limits_{t\to 0}\frac{D(t)}{t}$
        \end{itemize} 
        \item $f_{OUT}$ Output Flow Rate: speed of playout, considering 1x speed.
\end{itemize}
For a typical $f_{OUT} = 1$: $$\frac{dB}{dt} = \begin{cases}
            \frac{S}{R_c} - 1 \qquad \text{if playout going}\\
            \frac{S}{R_c} \qquad \text{if rebuffering, no playout}\\
        \end{cases}$$
The ideal condition is to have $S > R_c$ such that buffer is growing.\\